% Guide étudiant de StageCompagnon (mod_stage).
% Compilation : latexmk -pdf guide-etudiant.tex
\documentclass[11pt,a4paper]{article}

% Version du guide : à mettre à jour à chaque modification, avec une ligne dans l'historique.
\newcommand{\titreguide}{Guide étudiant}
\newcommand{\versionguide}{1.1}
\newcommand{\dateguide}{25 octobre 2026}
\newcommand{\versionstage}{0.5.1}
\newcommand{\versionsynthese}{0.2.1}

\usepackage{guide-commun}

\begin{document}

\pagedegarde{Suivre ses stages, demander une convention, s'auto-évaluer}{60a_student_top}

\tableofcontents

\section*{Historique des versions}
\addcontentsline{toc}{section}{Historique des versions}

\begin{center}
\begin{tabularx}{\linewidth}{@{}l l X@{}}
\toprule
\textbf{Version} & \textbf{Date} & \textbf{Modifications} \\
\midrule
1.0 & 16 octobre 2026 & Première version. \\
1.1 & 25 octobre 2026 & Captures d'écran refaites après le passage des listes et tableaux aux gabarits Mustache. \\
\bottomrule
\end{tabularx}
\end{center}

\newpage

% ===========================================================================
\section{Votre espace de stages}

Vos stages sont suivis dans le cours Moodle de votre promotion, par exemple
\og Stages -- Promotion 2030\fg{}, du début à la fin de votre cursus. Ouvrez
l'activité \og Stages\fg{} de ce cours : vous arrivez sur la page
\og Mes stages\fg{}.

\capture{60a_student_top}{Haut de la page \og Mes stages\fg{} : consignes, bouton
de demande de convention et synthèse.}

La page comprend :
\begin{itemize}
  \item la \textbf{synthèse} : durée retenue, années validées, thématiques
        obligatoires validées, mobilité internationale ;
  \item le \textbf{bilan par année d'étude}, jusqu'à votre année actuelle ;
  \item les \textbf{thématiques obligatoires}, avec la date limite et le bouton
        \bouton{Objectifs} ;
  \item la \textbf{mobilité internationale} exigée ;
  \item la liste de \textbf{tous vos stages}, avec leurs statuts et les actions
        possibles.
\end{itemize}

\capture{60b_student_bilan}{Bilan par année, thématiques obligatoires et mobilité.}

Seule la \textbf{durée retenue} compte : elle est fixée par la scolarité (DEVE)
à la validation du stage. Tant qu'un stage n'est pas validé, sa durée retenue
reste à~0.

% ===========================================================================
\section{Avant le stage}

\subsection{Consulter les objectifs}

Le bouton \bouton{Objectifs} d'une thématique ouvre ses documents à télécharger
et la check-list des objectifs à atteindre. Cette page est consultable à tout
moment, avant même d'avoir trouvé un stage.

\capture{62_objectives}{Objectifs d'une thématique.}

\subsection{Stage StageVet ou non}

\begin{itemize}
  \item \textbf{Chez un vétérinaire en France} : la convention est établie par
        StageVet. La scolarité enregistre le stage dans Moodle ; vous n'avez pas
        de demande à faire ici.
  \item \textbf{Autre stage} (structure non vétérinaire, étranger, etc.) : vous
        demandez la convention dans Moodle (§\ref{sec:demande}).
\end{itemize}

% ===========================================================================
\section{Demander une convention}
\label{sec:demande}

\begin{attention}
Déposez la demande au moins \textbf{2~semaines} avant le début du stage,
\textbf{4~semaines} pour un stage à l'étranger. Sans convention signée, le stage
ne peut pas commencer.
\end{attention}

\subsection{Deux points de départ}

\begin{itemize}
  \item Le stage n'apparaît pas dans votre liste :
        \bouton{Faire une demande de convention (hors StageVet)}. Un seul
        formulaire enregistre le stage et la demande.
  \item Le stage a été enregistré par la scolarité : \bouton{Demander la
        convention} sur sa ligne.
\end{itemize}

Il vous faut un enseignant référent. S'il n'en a pas été attribué, le formulaire
vous le signale : contactez la scolarité.

\subsection{Remplir le formulaire}

\capture{65_request_top}{Début du formulaire : stage, gabarit et référent.}

Tous les champs sont obligatoires, sauf le lieu du stage s'il est identique à
l'adresse de l'organisme. Le formulaire demande :
\begin{itemize}
  \item le stage : thématique, année d'étude, structure, durée, stage à
        l'étranger ;
  \item la langue et le modèle de convention ;
  \item votre enseignant référent, votre situation et le type de stage ;
  \item les plages de dates (plusieurs si le stage est en plusieurs fois) ;
  \item vos coordonnées, l'organisme d'accueil et le \textbf{tuteur}
        (maître de stage). Vérifiez son courriel : il recevra le lien
        d'évaluation ;
  \item les modalités particulières, la gratification, les congés ;
  \item la \textbf{check-list d'objectifs} de la thématique : cochez les
        objectifs que le stage permettra d'atteindre et justifiez chaque objectif
        non coché.
\end{itemize}

Cochez \bouton{Convention papier} si un exemplaire à signer à la main est
nécessaire.

\subsection{Suivi de la demande}

\begin{center}
\begin{tabularx}{\linewidth}{@{}l X@{}}
\toprule
\textbf{Statut de convention} & \textbf{Signification} \\
\midrule
\statut{Non demandée} & Aucune demande déposée. \\
\statut{En attente de l'enseignant référent} & Votre référent doit valider la demande. \\
\statut{Demandée} & La scolarité traite la demande. \\
\statut{Refusée} & À corriger : voir le motif reçu par courriel. \\
\statut{Éditée} & Convention prête : à télécharger, faire signer et rapporter. \\
\statut{Signée} & Convention signée : le stage peut avoir lieu. \\
\statut{Signée (StageVet)} & Convention établie par StageVet. \\
\statut{Sans convention} & Stage dispensé de convention. \\
\bottomrule
\end{tabularx}
\end{center}

\paragraph{Refus.} Vous recevez le motif par courriel, avec un lien vers votre
demande. Corrigez-la et soumettez-la à nouveau.

\paragraph{Convention prête.} Vous recevez un courriel. Sur la ligne du stage,
\bouton{Consulter la convention} télécharge le PDF. Faites-le signer par toutes
les parties et rapportez-le à la scolarité.

\capture{60c_student_stages}{Vos stages : un stage signé à auto-évaluer, une
convention éditée à télécharger, un stage dont la convention est à demander.}

Si la convention n'est pas signée 7~jours avant le début du stage, vous recevez
un rappel.

\begin{astuce}
Vous ne pouvez modifier un stage que tant qu'il est au statut
\statut{Enregistré} et que sa demande n'est pas en cours. Ensuite, adressez-vous
à la scolarité.
\end{astuce}

% ===========================================================================
\section{Pendant et après le stage}

\subsection{Évaluation par le maître de stage}

À partir du premier jour du stage, votre maître de stage reçoit par courriel un
lien vers un questionnaire, sans création de compte. Prévenez-le.

\subsection{Votre auto-évaluation}

Le bouton \bouton{Auto-évaluer mon stage} est utilisable dès le premier jour du
stage, une fois la convention signée.

\capture{64_selfeval}{Auto-évaluation : jours effectifs et questionnaire.}

\begin{enumerate}
  \item \textbf{Jours de stage effectifs} : cochez les jours réellement
        travaillés, puis \bouton{Enregistrer}. La réglementation impose au moins
        un jour de repos par semaine.
  \item \textbf{Rapport de stage}, si la thématique le demande : déposez vos
        documents.
  \item \textbf{Questionnaire} : répondez aux questions (les questions marquées
        d'un astérisque sont obligatoires), puis \bouton{Enregistrer}.
\end{enumerate}

\begin{attention}
Une auto-évaluation enregistrée n'est plus modifiable. En cas d'erreur, demandez
à la scolarité de rouvrir la saisie.
\end{attention}

Le stage passe à \statut{Évalué par l'étudiant} et votre référent est prévenu.

\subsection{Validation}

\begin{enumerate}
  \item Votre référent évalue le stage : \statut{Évalué par l'enseignant}.
  \item La scolarité fixe la durée retenue et valide : \statut{Validé DEVE}.
        Votre bilan est mis à jour.
\end{enumerate}

Un stage \statut{Non validé} ou \statut{Annulé} ne compte pas : contactez la
scolarité.

% ===========================================================================
\section{Courriels que vous recevez}

\begin{center}
\begin{tabularx}{\linewidth}{@{}l X@{}}
\toprule
\textbf{Courriel} & \textbf{Ce qu'il faut faire} \\
\midrule
Convention refusée & Corriger la demande et la soumettre de nouveau. \\
Convention disponible & Télécharger, faire signer, rapporter à la scolarité. \\
Convention non signée & Le stage commence dans moins de 7~jours : régulariser au plus vite. \\
\bottomrule
\end{tabularx}
\end{center}

Les courriels partent d'une adresse \og no-reply\fg{} : n'y répondez pas,
écrivez à la scolarité.

% ===========================================================================
\section{En résumé}

\begin{enumerate}
  \item Consulter les objectifs de la thématique.
  \item Stage hors StageVet : demander la convention 2~semaines avant
        (4~à l'étranger).
  \item Télécharger la convention éditée, la faire signer, la rapporter.
  \item Prévenir le maître de stage qu'il recevra un lien d'évaluation.
  \item Dès le premier jour : renseigner les jours effectifs, puis
        s'auto-évaluer en fin de stage.
  \item Vérifier dans \og Mes stages\fg{} que le stage est validé.
\end{enumerate}

\end{document}
